\section{Positive sectors of the Weil form}
\label{intake:weil:section}

The broad dilations in Proposition~\ref{prop:gamma-negative}
concentrate the Fourier mass near zero. Narrow support instead forces
enough Fourier mass into the region where the archimedean kernel is positive.
The prime-power terms of the full Weil form couple separated intervals.
A sufficiently small width makes the archimedean contribution dominate these couplings.
This gives a positive space of functions with infinitely many independent directions,
while retaining the full arithmetic functional.

\subsection{Conventions and the full form}
For $f\in C_c^\infty(\mathbb R)$ use
\[
 \widehat f(t)=\int_\mathbb R f(u)e^{itu}\,du,\qquad
 f^*(u)=\overline{f(-u)},\qquad
 (f*g)(u)=\int_\mathbb R f(v)g(u-v)\,dv.
\]
Let $\vmLambda(n)$ be the von Mangoldt function: it is $\log p$ for $n=p^a$ and zero otherwise. Put
\[
 k_\infty(t)=\operatorname{Re}\psiGamma(1/4+it/2)-\log\pi,
\]
where $\psiGamma=\Gamma'/\Gamma$. Define the full arithmetic functional by
\[
\begin{split}
 L_\zeta(h)={}&\widehat h(i/2)+\widehat h(-i/2)
 +\frac1{2\pi}\int_\mathbb R\widehat h(t)k_\infty(t)\,dt\\
 &-\sum_{n\ge2}\frac{\vmLambda(n)}{\sqrt n}
          \bigl(h(\log n)+h(-\log n)\bigr).
\end{split}
\]
This is the full functional in Theorem~\ref{thm:weil}, restricted to
compact smooth tests, with exactly the same Fourier normalization.
For compact smooth $h$ the prime sum is finite and the Fourier integral converges. Define
\[
 B_\zeta(f,g)=L_\zeta(f^**g),\qquad Q_\zeta(f)=B_\zeta(f,f).
\]
The form is conjugate-linear in its first argument, linear in its second and Hermitian. The estimates below use this displayed functional directly.

\subsection{A coercive estimate on one narrow interval}
The convergent digamma series gives
\[
 k_\infty(t)=k_\infty(0)+\sum_{n\ge0}
 \left(\frac1{n+1/4}
 -\frac{n+1/4}{(n+1/4)^2+t^2/4}\right).
\]
Every summand is nonnegative and increasing with $|t|$. Also
\[
 k_\infty(0)=-\gamma-\frac\pi2-3\log2-\log\pi>-6.
\]
For example $\gamma<1$, $\pi<22/7$, $\log2<7/10$ and $\log\pi<6/5$ suffice. The last two exponential inequalities follow from their finite positive Taylor lower bounds. Finally $k_\infty(T)\to+\infty$: for $n+1/4\le T/2$, the corresponding summand in the displayed digamma series is at least $1/(2(n+1/4))$, whose partial sums diverge.

Let $f$ have support in an interval of length $L$. Cauchy--Schwarz gives $|\widehat f(t)|^2\le L\|f\|_2^2$, and Plancherel gives
\[
 \frac1{2\pi}\int_{-T}^T|\widehat f(t)|^2dt
       \le\frac{LT}{\pi}\|f\|_2^2,\qquad
 \frac1{2\pi}\int_\mathbb R|\widehat f(t)|^2dt=\|f\|_2^2.
\]
Let $T>0$ and $M\ge-6$ satisfy $k_\infty(T)\ge M$.
The archimedean contribution is then at least
\[
 \left(M-(M+6)\frac{LT}{\pi}\right)\|f\|_2^2.
\]
When $L<\log2$, the support of $f^**f$ contains no nonzero prime-power logarithm, so its prime contribution is zero. The polar term is bounded below by $-2Le^{L/2}\|f\|_2^2$: translate the interval to be centred at zero and apply Cauchy--Schwarz to $\int f(u)e^{u/2}du$ and $\int f(u)e^{-u/2}du$. The translation factors cancel between these two integrals. Consequently
\[
 Q_\zeta(f)\ge c(L,T,M)\|f\|_2^2,
 \qquad c(L,T,M)=M-(M+6)\frac{LT}{\pi}-2Le^{L/2}.
\]
The estimate holds uniformly over all smooth functions in the interval.

\subsection{The cross-window kernel retains the arithmetic primes}
Away from zero, the distribution given by the sum of the polar and archimedean contributions in the displayed definition of $L_\zeta$ is the smooth even kernel
\[
 K(u)=2\cosh(u/2)-\frac{e^{-|u|/2}}{1-e^{-2|u|}}\qquad(u\ne0).
\]
To verify it, use the integral expression
\[
 \psiGamma(a+it/2)=-\gamma+
 \int_0^\infty\frac{e^{-v}-e^{-(a+it/2)v}}{1-e^{-v}}\,dv.
\]
Test on a smooth function supported away from zero. Terms constant in $t$ vanish on that test. Fourier inversion of the cosine term, followed by $v=2|u|$, gives the negative second summand in the displayed formula for $K$. The two polar evaluations give $2\cosh(u/2)$. This argument is on separated supports; it does not ignore a delta term at zero in the diagonal form.

Fix distinct positions $a_1,\ldots,a_N$. Let $I_L=(-L/2,L/2)$ and
\[
 F(u)=\sum_{j=1}^N f_j(u-a_j),\qquad f_j\in C_c^\infty(I_L).
\]
Choose $L$ so small that these supports are disjoint and each difference interval contains a prime-power logarithm only if its centre equals that logarithm. This is possible because there are only finitely many prime powers in a bounded interval. Put
\[
 w_{ij}=\begin{cases}
 \vmLambda(n)/\sqrt n,&|a_i-a_j|=\log n\text{ for an integer }n\ge2,\\
 0,&\text{otherwise}.
 \end{cases}
\]
For a resonant pair its prime term is exactly $-w_{ij}\int f_j(v)\overline{f_i(v)}dv$. The smooth part is the kernel $K$ evaluated at $a_j-a_i+v-u$. If $M_{ij}$ bounds $|K|$ there, its absolute value is at most $LM_{ij}\|f_i\|_2\|f_j\|_2$. Thus
\[
 Q_\zeta(F)\ge
 \left(c(L,T,M)-\max_i\sum_{j\ne i}(w_{ij}+LM_{ij})\right)
       \sum_i\|f_i\|_2^2.
\]
Apply $2ab\le a^2+b^2$ to each unordered pair to obtain this bound.

\begin{theorem}[Positive arithmetic windows of arbitrary finite multiplicity]\label{intake:weil:thm:any-windows}
For any fixed finite set of distinct real positions $a_1,\ldots,a_N$, there exists $L>0$ such that the full Weil form is positive definite on the infinite-dimensional space
\[
 \bigoplus_{j=1}^N C_c^\infty(a_j+I_L).
\]
It has a uniform strictly positive lower bound by the $L^2$ norm. All prime-power terms in the displayed definition of $L_\zeta$ are retained.
\end{theorem}
\begin{proof}
The finite row sums of the $w_{ij}$ are fixed. Choose $M$ larger than their maximum plus two, and then $T$ with $k_\infty(T)\ge M$. For sufficiently small $L$, the separation requirements hold, $LT$ and $Le^{L/2}$ tend to zero, and the $M_{ij}$ remain bounded because their arguments stay away from zero. The coefficient in the row-sum estimate is then positive, in fact larger than one after making $L$ smaller. 
\end{proof}

\subsection{An explicit four-window theorem with a prime-power coupling}
Choose
\[
 (a_0,a_1,a_2,a_3)=(0,\log2,\log3,\log4),\qquad L=10^{-3}.
\]
The resonant pairs are $(0,1),(0,2),(0,3),(1,3)$; their weights are respectively
\[
 \frac{\log2}{\sqrt2},\quad\frac{\log3}{\sqrt3},\quad
 \frac{\log2}{2},\quad\frac{\log2}{\sqrt2}.
\]
In particular the $\log4$ term is a genuine prime-power term, not $\log4/2$. Every other logarithmic difference is separated from prime-power logarithms by more than $L$.

\begin{theorem}[A quantitative positive sector]\label{intake:weil:thm:four-windows}
For arbitrary $f_j\in C_c^\infty((-1/2000,1/2000))$, the function
\[
 F(u)=f_0(u)+f_1(u-\log2)+f_2(u-\log3)+f_3(u-\log4)
\]
satisfies
\[
 \boxed{\displaystyle Q_\zeta(F)\ge\frac15\sum_{j=0}^3\|f_j\|_2^2.}
\]
\end{theorem}
\begin{proof}
A finite rational calculation gives
\[
 \sum_{n=0}^{99}\left(\frac1{n+1/4}
 -\frac{n+1/4}{(n+1/4)^2+2500}\right)>8.
\]
For an explicit integer calculation, put
\[
 d_n=(4n+1)((4n+1)^2+40000),\qquad
 b_j=\sum_{n=10j}^{10j+9}\left\lfloor\frac{1600000000}{d_n}\right\rfloor.
\]
The values for $j=0,\ldots,9$ are
\[
 (b_0,\ldots,b_9)=(64860,6514,3300,1958,1248,832,574,411,300,226).
\]
Their sum is $80223$. The rational partial sum is therefore at least
$80223/10000>8$. Since all later summands are nonnegative and $k_\infty(0)>-6$, we have $k_\infty(100)>2$. The local coercivity estimate, $\pi>3$, and $e^{1/2000}<3/2$ give a diagonal lower bound
\[
 2-\frac8{30}-\frac3{1000}=\frac{5191}{3000}.
\]
All cross-kernel arguments have absolute value between $\log(4/3)-10^{-3}$ and $\log4+10^{-3}$. On this interval each of the two positive terms in the displayed formula for $K$ is less than three. For the denominator term, $\log(4/3)>1/4$ and $\log(3/2)<1/2-1/8+1/24<0.498$ imply $1-e^{-2|u|}>1/3$. For the polar term use $2\cosh((\log4+10^{-3})/2)<2(1.001)+1/2<3$. Thus $|K(u)|<3$, and the smooth row cost is less than $9/1000$.

The largest prime row sum is bounded by
\[
 \frac12+\frac{11}{17}+\frac7{20}=\frac{509}{340},
\]
using $\log2<7/10$, $\sqrt2>7/5$, $\log3<11/10$, $\sqrt3>17/10$. Substitution into the row-sum estimate gives
\[
 \frac{5191}{3000}-\frac9{1000}-\frac{509}{340}=\frac{5719}{25500}>\frac15.
\]
This proves the stated four-window inequality for the entire stated infinite-dimensional space.
\end{proof}

This sector is not a convolution algebra. Its support windows are not preserved under repeated convolution, and the permissible window width in Theorem~\ref{intake:weil:thm:any-windows} depends on the whole finite set of positions. Neither theorem supplies a dense convolution-stable exhaustion with one uniform estimate. That missing extension is exactly where the global Weil-positivity problem remains. Global positivity on the full convolution algebra is not a consequence of these estimates.


\subsection{Finite-rank operators from the arithmetic pairing}

Let $\mathscr H$ be the Hilbert completion of this four-window space for
$B_\zeta$, with the inner product conjugate-linear in its first argument.
For $f,g\in\mathscr H$, define the bounded rank-one operator
$\Theta_{f,g}h=B_\zeta(g,h)f$. Direct substitution gives
\[
 \Theta_{f,g}\Theta_{u,v}=B_\zeta(g,u)\Theta_{f,v},\qquad
 \Theta_{f,g}^*=\Theta_{g,f},\qquad
 \operatorname{Tr}(\Theta_{f,g}\Theta_{u,v})
 =B_\zeta(g,u)B_\zeta(v,f).
\]
The trace identity follows from
$\operatorname{Tr}(\Theta_{f,v})=B_\zeta(v,f)$ in an orthonormal basis.
For vectors in the original test space, every displayed contraction has
its prime, archimedean and polar terms given by the full functional above.
These operators represent the selected quadratic pairing. A representation
of convolution would additionally require compatible operators for every
convolution multiplier on a common invariant domain.


For a finite-dimensional flavour subspace $V\subset\mathscr H$, take the
charged fermion space
\[
 \mathcal F_V=\Lambda_{\mathbb C}\!\left(
 \bigoplus_{r\in\mathbb Z_{>0}-1/2}(V\oplus\overline V)t^r\right).
\]
Use its algebraic finite-particle domain and the induced positive Hermitian form.
Exterior multiplication and its adjoint satisfy the canonical
anticommutation relations with contraction $B_\zeta$.
Writing the resulting charged fields as $\psi_f(z)$, linear in $f$,
and $\psi_g^*(w)$, conjugate-linear in $g$,
the vacuum contraction in the expansion $|w|<|z|$ is
\[
 B_\zeta(g,f)\sum_{k\ge0}z^{-k-1}w^k
 =\frac{B_\zeta(g,f)}{z-w}.
\]
For a finite-rank $A$, choose $V$ containing the ranges of $A$ and $A^*$,
and define $J_{\Theta_{f,g}}(z)=:\psi_f(z)\psi_g^*(z):$, extending linearly in $A$.
Normal ordering places creation operators to the left, with fermion signs.
The singular operator product is
\[
 J_A(z)J_B(w)\sim
 \frac{\operatorname{Tr}(AB)}{(z-w)^2}
 +\frac{J_{[A,B]}(w)}{z-w}.
\]
Indeed the two single contractions compose the matrices in their two
orders, with opposite fermion signs. The double contraction sums their
closed index pair and gives $\operatorname{Tr}(AB)$.
Equivalently, the current modes satisfy
\[
 [J_A(m),J_B(n)]=J_{[A,B]}(m+n)
 +m\delta_{m+n,0}\operatorname{Tr}(AB)I.
\]
For $m>0$, exactly $m$ half-integer modes cross the vacuum in the scalar
term. Antisymmetry determines the remaining signs.
For fixed mode and finite-energy vector, only finitely many creation
terms survive. Enlarging $V$ preserves these identities, so they apply
to every finite set of finite-rank operators on $\mathscr H$.
The arithmetic pairing thus occurs as the contraction of singular
currents on this specified positive sector. Its support is still not
stable under convolution.
